% =====================================================================
\section{Literature and Technical Review}
% =====================================================================

\subsection{Overview of Existing Solutions}

Three groups of solutions are relevant. Table~\ref{tab:existing} summarizes how each one relates to
the requirements of this project.

\textbf{Commercial ALPR services and cloud OCR.} Rekor (the company behind OpenALPR)
\cite{openalpr}, Plate Recognizer \cite{platerecognizer} and the text detection feature of Amazon
Rekognition \cite{awsrekognition} are mature products, and dedicated ALPR products are common in
parking access control. They handle many countries and camera types, scale on demand and come with
documentation and support. They are built around one image per
request, and the cloud versions process images on the provider's servers, with a price per call.

\textbf{Open-source OCR libraries.} The open-source OpenALPR engine \cite{openalpr} and general OCR
systems such as PP-OCR \cite{du2020ppocr} can be self-hosted and adapted. They are designed for
images in which the characters are legible, and they read one image at a time.

\textbf{Research models for low-resolution plates.} One line of work restores the image before
reading it, using attention and sub-pixel convolution \cite{nascimento2023attention} or layout- and
character-aware losses \cite{nascimento2024lcdnet}. Another line combines several frames: voting over
super-resolved frames \cite{nascimento2025toward} and optical-flow restoration in MF-LPR$^2$
\cite{na2025mflpr2}. Most recognizers in this area build on CRNN with CTC
\cite{shi2017crnn,graves2006ctc}. In the ICPR 2026 competition, the best entry reached $82.13\%$ on
the competition test set \cite{laroca2026icpr}. These models are released as training and evaluation
code, which suits research use.

\textbf{Selective prediction.} Selective classification \cite{geifman2017selective} lets a model
abstain when its confidence is low and studies the trade-off between coverage (the share of inputs
answered) and the accuracy on the answered inputs. This is the basis for the review queue.

\begin{table}[tbp]
\centering
\caption{Existing solutions against the needs of this project.}
\label{tab:existing}
\begin{tabularx}{\textwidth}{@{}P{25mm} P{17mm} P{19mm} L@{}}
\toprule
Solution group & Input & Deployment & Fit with this project \\
\midrule
Commercial ALPR, cloud OCR & one image & cloud or licensed SDK & Designed for legible plates; images leave the site, which the dataset license does not allow \\
Open-source OCR & one image & self-hosted & Can run on site; designed for legible characters \\
Research LR models & 1--5 frames & training code & Target the right images; packaged for experiments rather than for operators \\
\textbf{PlateReader LR} & 1--5 frames & self-hosted service & Adds API, review workflow, storage and deployment; limited to Brazilian layouts \\
\bottomrule
\end{tabularx}
\end{table}

\subsection{Technical Gaps for This Use Case}

The solutions above serve their intended settings. For low-resolution surveillance tracks, five
needs remain open in the material reviewed.

\begin{description}
  \item[G1. Resolution.] At five to seven pixels per character, the input is below the range that
        general ALPR and OCR engines document, and closed engines cannot be retrained on it.
  \item[G2. Multiple frames.] Single-image interfaces cannot use the extra evidence in the other
        frames of a track.
  \item[G3. Abstention and review.] The tools reviewed return a reading for every input. None of
        them routes low-confidence readings to a reviewer with aids for correction.
  \item[G4. Data and cost constraints.] The dataset license allows academic use only and forbids
        redistribution \cite{icpr26license}, so images cannot be sent to a third-party cloud. The
        target sites have CPUs but no GPU.
  \item[G5. Engineering.] Research models usually come without an API, persistent storage, model
        versioning or a repeatable evaluation pipeline, which are needed to run a model as a
        service \cite{sculley2015hidden}.
\end{description}

\subsection{Rationale for the Selected Approach and Stack}

\textbf{Model.} MF-SR-OCR addresses G1 and G2. It rectifies, aligns, super-resolves and fuses five
frames in one network and reaches $79.95\%\pm0.15$ over three seeds, against $75.78\%$ for the
retrained CRNN+STN baseline (Fig.~\ref{fig:multiseed}). A variant without the super-resolution
branch (J1) reaches $80.45\%\pm0.45$; paired McNemar tests per seed give $p=0.91$, $0.43$ and
$0.64$, so the two are statistically indistinguishable on this data \cite{nguyen2026mfsrocr}.
MF-SR-OCR is served because it has the lowest seed-to-seed spread and because its super-resolved
frames can be shown to reviewers. Whether those frames help reviewers is tested in RQ3. J1 needs
$4.17\times$ fewer operations per track and remains available as a fallback through configuration.

\textbf{Architecture.} Table~\ref{tab:arch} compares three options. A modular monolith fits one
developer, one CPU host and one model that should be loaded once.

\begin{table}[tbp]
\centering
\caption{Architecture options.}
\label{tab:arch}
\begin{tabularx}{\textwidth}{@{}P{22mm} L L@{}}
\toprule
Option & Strengths & Costs for this project \\
\midrule
\textbf{Modular monolith} (chosen) & One process loads the 29.6M-parameter model once; simple to test and deploy; module boundaries kept in code \cite{fowler_monolithfirst} & Scales as one unit; a split is possible later along the existing modules \\
Microservices & Independent scaling and deployment & Network calls between services, several images to run and monitor; more than one developer can maintain in 12 weeks \\
Serverless functions & No server to manage & Cold starts would reload the model; GPU-free functions with a 120\,MB model are slow to start; images would leave the site \\
\bottomrule
\end{tabularx}
\end{table}

\textbf{Database.} The workload mixes long batch writes with frequent reads from the review screen.
Table~\ref{tab:db} compares the options. PostgreSQL was chosen; in a Week~2 load test, review-queue
reads stayed at $46$\,ms p95 with no errors while a batch was writing (Section~\ref{sec:vv}).

\begin{table}[tbp]
\centering
\caption{Database options for the result and review store.}
\label{tab:db}
\begin{tabularx}{\textwidth}{@{}P{22mm} L L@{}}
\toprule
Option & Strengths & Costs for this project \\
\midrule
SQLite (first roadmap) & No server; one file & One writer at a time \cite{sqlite_lock}, so reviewers wait behind batch writes; no hot backup tool \\
\textbf{PostgreSQL 16} (chosen) & MVCC: reads and writes do not block each other \cite{postgres_mvcc}; \code{JSONB}, \code{UUID}, \code{CHECK}; hot backup with \code{pg\_dump} & One more container to run and back up \\
MongoDB & Flexible documents & The data is relational (batch, prediction, review); joins and constraints would move into application code \\
\bottomrule
\end{tabularx}
\end{table}

\textbf{Backend, API style and frontend.} FastAPI \cite{fastapi} keeps the service in Python, so the
API reuses the exact preprocessing of the training code; this is what makes the 999/999 parity
possible. It also generates the OpenAPI description. REST with JSON fits browser clients and file
uploads; errors follow RFC~9457 \cite{rfc9457}. gRPC would add a proxy layer for browsers without a
clear gain at this load. Batch progress uses polling for now; Server-Sent Events are an option if
polling becomes a cost. For the frontend, Gradio, Streamlit and React \cite{react} were compared in
Week~1. React was chosen because the review screen needs custom interaction (suggestion buttons,
per-character highlights) and the author has production experience with it. Its extra cost of two to
three weeks is contained by freezing the interface at three screens after Week~8.

\textbf{Deployment.} Docker Compose runs the whole stack on one host with one command; Kubernetes
would add a control plane that a single CPU host does not need. Cloudflare Tunnel with Cloudflare
Access \cite{cloudflare_tunnel} gives a public demo address without opening an inbound port and adds
sign-in in front of an API that has no accounts of its own. Table~\ref{tab:stack} summarizes the
stack.

\begin{table}[tbp]
\centering
\caption{Technology stack.}
\label{tab:stack}
\begin{tabularx}{\textwidth}{@{}>{\bfseries}P{20mm} P{34mm} L@{}}
\toprule
Layer & Choice & Note \\
\midrule
Model & PyTorch, MF-SR-OCR (EMA weights) & CPU inference, 0.19--0.27\,s per track natively \\
Backend & Python 3.11, FastAPI, Uvicorn & Pydantic schemas form the API contract \\
Database & PostgreSQL 16, psycopg 3 pool & Text results only; no images \\
Frontend & React, Vite, Tailwind CSS & Three screens, desktop and mobile \\
Packaging & Docker, Docker Compose & Profiles for optional services \\
Access & Cloudflare Tunnel and Access & Outbound-only connection, sign-in at the edge \\
CI/CD & GitHub Actions \cite{github_actions}, GHCR & Lint, tests, scans, image build \\
\bottomrule
\end{tabularx}
\end{table}
